% Generated draft: requires mathematical and source review.
\documentclass[11pt]{article}
\usepackage[margin=25mm]{geometry}
\usepackage{fontspec}
\setmainfont{Latin Modern Roman}
\usepackage{amsmath,amssymb,mathtools,unicode-math}
\setmathfont{Latin Modern Math}
\usepackage{enumitem,needspace,xurl,hyperref}
\setlength{\parindent}{0pt}
\setlength{\parskip}{5pt}
\setlength{\emergencystretch}{3em}
\newcommand{\N}{\mathbb{N}}
\newcommand{\Z}{\mathbb{Z}}
\newcommand{\Q}{\mathbb{Q}}
\newcommand{\R}{\mathbb{R}}
\newcommand{\C}{\mathbb{C}}
\newcommand{\F}{\mathbb{F}}
\newcommand{\A}{\mathbb{A}}
\newcommand{\PP}{\mathbb{P}}
\newcommand{\E}{\mathbb{E}}
\newcommand{\eps}{\varepsilon}
\newcommand{\dd}{\,\mathrm{d}}
\begin{document}
\section*{UnsolvedMath ID: 800362}
\textbf{Source problem code:} catalogue entry
\subsection{Definitions and mathematical statement}
Let \((W,S)\) be a finite Weyl group of type \(B_n\) or \(D_n\), and let \(\ell:W\to\mathbb Z_{\geq 0}\) be its Coxeter length function. The source assumes the equal-parameter condition \(L=\ell\). It does not specify the lower-rank conventions for these types.

The supplied excerpt does not define the coefficient ring \(A\) or the normalization of the Iwahori--Hecke algebra \(H\). The following explicit formulation is conditional on the normalized equal-parameter convention: set \(A=\mathbb Z[v,v^{-1}]\), with \(v\) an indeterminate, and let \(H\) be the free \(A\)-module with basis \(\{T_x:x\in W\}\), multiplication determined by
\[
T_xT_y=T_{xy}\quad\text{if }\ell(xy)=\ell(x)+\ell(y),
\qquad
(T_s-v)(T_s+v^{-1})=0\quad(s\in S).
\]
This convention must be checked against the original source before the draft is accepted.

For a conjugacy class \(C\subseteq W\), define
\[
\ell(C)=\min\{\ell(x):x\in C\},
\qquad
C_{\min}=\{x\in C:\ell(x)=\ell(C)\}.
\]
For \(w\in W\), define the \(A\)-linear endomorphism
\[
\Phi_w:H\longrightarrow H,
\qquad
h\longmapsto v^{2\ell(w)}T_whT_{w^{-1}},
\]
and set
\[
N^w=\operatorname{Tr}_A(\Phi_w).
\]
Here the extracted expression \(T_{w\backslash\mathrm{i}}\) is interpreted as \(T_{w^{-1}}\), the basis element indexed by the group inverse; it is not interpreted as the algebra inverse \(T_w^{-1}\). The extraction does not itself resolve this notation.

The source states that \(N^w\in\mathbb Z[v^2]\), that
\[
\left.N^w\right|_{v=1}=|C_W(w)|,
\qquad
C_W(w)=\{x\in W:xw=wx\},
\]
and that, by a consequence of its reference [GP], \(N^w\) is independent of the choice of \(w\in C_{\min}\). Using this asserted independence, write \(N_C=N^w\) for \(w\in C_{\min}\). Call \(C\) positive if
\[
N_C\in\mathbb Z_{\geq 0}[v^2],
\]
meaning that every coefficient is nonnegative. The minimum-length restriction is essential to this definition.

\textbf{Problem.} For every finite Weyl group \(W\) of type \(B_n\) or \(D_n\), give a complete list of its positive conjugacy classes.

As background, the supplied source states that elliptic regular conjugacy classes, in the sense of [S74], are positive, citing [L18]; that in type \(A_n\) the positive classes are the identity class and the Coxeter-element class; and that [L18] gives a complete list in exceptional Weyl types. These are supplied background assertions, not independently verified claims. The requested classification remains conditional on confirming the algebra conventions and the extracted inverse notation.
\subsection{Short English statement}
For each finite Weyl group of type \(B_n\) or \(D_n\), classify the conjugacy classes whose specified Hecke-algebra trace polynomial has only nonnegative coefficients, using a minimum-length representative of each class. The supplied record labels the problem open. The algebra normalization, inverse notation, and lower-rank conventions require confirmation.
\subsection{Sources}
\begin{itemize}
\item Local UnsolvedMath registry, permanent ID \(800362\), source problem code \texttt{catalogue entry}, title ``Iwahori--Hecke Algebras --- Problem 3,'' in the AMR Open Problem Lists collection. The supplied record attributes the problem to Lusztig (2020) and labels it \texttt{open}; this label is supplied metadata, not an independently verified status assessment.
\item Supplied primary-source attribution: Lusztig, \emph{Open Problems on Iwahori--Hecke Algebras} (2020), Problem 3, part (d), \url{https://arxiv.org/abs/2006.08535}. The background records extraction by \texttt{arxiv\_latex} and an access date of 2026-10-16. That date is supplied provenance, not a claim that this draft independently accessed or verified the source.
\item The supplied excerpt cites [GP] for the deduction that \(N^w\) is constant on \(C_{\min}\), [S74] for the meaning of elliptic regular conjugacy classes, and [L18] for their positivity and the exceptional-type lists. Full bibliographic entries for these labels are not supplied.
\item The extraction ends with a formatting fragment containing the key \texttt{L78}. Its bibliographic meaning and connection to the requested classification are not supplied.
\end{itemize}
\end{document}
